\documentclass[11pt]{article}
\usepackage[a4paper,margin=22mm]{geometry}
\usepackage{amsmath,amssymb}
\usepackage{graphicx}
\usepackage{booktabs,longtable}
\usepackage{caption}
\usepackage{url}
\usepackage[hidelinks]{hyperref}
\captionsetup{font=small,labelfont=bf,justification=justified,singlelinecheck=false}
\setlength{\parskip}{0.4em}
\setlength{\parindent}{0pt}
\setlength{\emergencystretch}{3em}
\newcommand{\R}{\mathbb{R}}

\title{Regional structure of shared neural variance in IBL mesoscope recordings}
\author{Claude \and International Brain Laboratory (IBL) \and Michael Schartner}
\date{\today}

\begin{document}
\maketitle

\section{Methods}

\subsection{Data}

All analyses use two-photon mesoscope recordings from the International Brain Laboratory (IBL; lab \texttt{cortexlab}, subjects SP037--SP081), taken from the curated list of 165 ``canonical'' sessions in \texttt{canonical\_sessions.txt}. Per field of view (FOV) we load the deconvolved activity (\texttt{mpci.ROIActivityDeconvolved}), the neuron/non-neuron classification of each region of interest (ROI; \texttt{mpciROIs.mpciROITypes}), the ROI's Allen CCF~2017 brain-region id (histology-registered \texttt{brainLocationIds\_ccf\_2017}, falling back to the registration-pipeline estimate when histology is not yet available) and its estimated medio-lateral (ML), antero-posterior (AP) and dorso-ventral coordinates in \textmu m (\texttt{mpciROIs.mlapdv\_estimate}). Each ROI's timestamps are the imaging frame times plus the ROI's scan-line-dependent sub-frame offset (\texttt{mpciStack.timeshift}). Only ROIs classified as neurons are kept, and all FOVs of a session are stacked.

A neuron belongs to a region if its Allen acronym starts with the region's acronym, so that all layers are pooled: primary visual cortex (VISp), barrel-field primary somatosensory cortex (SSp-bfd), primary motor cortex (MOp) and secondary motor cortex (MOs).

Behavioural and task variables come from the same sessions' IBL extractions: wheel position and timestamps, the trials table (block prior \texttt{probabilityLeft}, \texttt{choice}, \texttt{firstMovement\_times}) and the raw left-camera video with its frame timestamps.

\subsection{Session selection}

To find eligible sessions, all canonical sessions were scanned for the number of neurons in each region, the imaged duration and the availability of trials and video. Only per-FOV metadata were used for this scan, not activity, with the same neuron and region definitions as the analyses. 161 of the 165 sessions could be scanned; the remaining four have inconsistent ROI array sizes and cannot be loaded.

\paragraph{Dimensionality analysis (Fig.~\ref{fig:dim}).} For each region, the five sessions with the most neurons in that region were used, among sessions with at least 1,000 neurons there (Table~\ref{tab:dim}). This gives 20 region-sessions: VISp five sessions from four animals, SSp-bfd five from four, MOp five from two and MOs five from one animal (SP076, the only animal with MOs recordings).

\paragraph{Behavioural prediction analysis (Fig.~\ref{fig:behav}).} Sessions had to have at least 1,000 neurons in the region, an extracted trials table, a raw left-camera video, and imaging covering the analysis window (1,730--2,030~s). Each region has eight sessions (Table~\ref{tab:behav}): the three sessions of the earlier version of this figure, plus five sessions chosen to maximise the number of animals, taking sessions from animals not yet represented first and, within that, the session with the most neurons in the region. This gives VISp eight sessions from eight animals, SSp-bfd eight from eight, MOp eight from three and MOs eight from one animal (SP076). Four sessions contribute to two regions (three MOp/MOs, one VISp/SSp-bfd), always with disjoint neuron sets.

\paragraph{Passive-protocol analyses (Figs.~\ref{fig:pdim} and \ref{fig:pbehav}).} In most sessions the task is followed by a passive protocol, whose periods are given by the extracted \texttt{passivePeriods} table: spontaneous activity in front of a grey screen, receptive-field mapping with sparse noise, and a replay of task stimuli without the wheel being coupled to the stimulus. Because this epoch is shorter than the task, the passive analyses use the first 730~s of imaged passive protocol, starting at the first imaging frame after the protocol began (the longest window available in nearly all sessions). Depending on the session this window contains 60, 300 or 600~s of spontaneous activity, followed by receptive-field mapping and, if time remains, the start of task replay. The window is therefore not stimulus-free: in the 13 of 20 sessions with a 60-s spontaneous block, receptive-field mapping ($\sim$300~s) and task replay ($\sim$330~s) make up most of it. A stimulus-free window would be far too short: the spontaneous block alone lasts only 60~s in most sessions, and spontaneous activity plus receptive-field mapping ($\sim$370~s) is about as short as a 400-s window, which gave unreliable spectra ($\sim$30\% of ranks 10--100 had $r_k \le 0$). The 730-s window is thus a compromise between the length needed for stable estimates and stimulus content. Sessions had to have at least 1,000 neurons in the region and at least 730~s of imaged passive protocol. Of the 20 dimensionality sessions, three have no passive protocol and two have too little imaged passive data. The remaining 15 were kept and five replacements were added, again preferring animals not yet represented, giving five sessions per region (VISp four animals, SSp-bfd four, MOp three, MOs one). For the passive behavioural analysis, the two MOs sessions without camera data were replaced by the next SP076 sessions with video (Table~\ref{tab:passive}).

\subsection{Making sessions and regions comparable}

Every region-session is analysed independently and identically; neural activity is never concatenated across sessions or animals, and each session contributes one curve with equal weight. Specifically:
\begin{enumerate}
  \item \textbf{Identical neuron count.} From each region-session exactly $N = 1{,}000$ neurons are drawn uniformly at random without replacement from all neurons of that region (fixed seed 0), so that neither the number of neurons nor its effect on the noise level can differ between regions or sessions. After the spatial split (below) the two halves contain 466--534 neurons each.
  \item \textbf{Identical time window.} All sessions of an analysis use the same interval on the session clock: 100--3,100~s (3,000~s) for dimensionality and 1,730--2,030~s (300~s) for behavioural prediction. The passive analyses instead use a window of identical duration (730~s) aligned to the start of each session's passive protocol, and the passive dimensionality analysis is repeated on a duration-matched task window (100--830~s) from the same neurons, so that passive--task differences cannot arise from the shorter recording.
  \item \textbf{Identical processing.} The same temporal binning, spatial split rule, train/test partition, number of components ($K_{\max}=128$), predictor definitions and random seeds are used for all region-sessions.
\end{enumerate}
Two differences remain that could not be matched. (i) The imaging frame period differs between sessions (0.18--0.29~s, depending on the number of FOVs and planes; Table~\ref{tab:dim}). For the dimensionality analysis this is equalised by averaging into 1.25-s bins; the behavioural analysis is run at the native frame rate, so its sessions have different numbers of time points (964--1,929 frames in the 300-s window; Table~\ref{tab:behav}). (ii) In five dimensionality sessions (four VISp, one SSp-bfd) the task ended 35--230~s before 3,100~s, so the end of the window contains post-task time (partly unimaged, partly the start of the passive protocol; Table~\ref{tab:dim}).

\subsection{Notation}

\begin{longtable}{@{}p{0.2\textwidth}p{0.75\textwidth}@{}}
\toprule
Symbol & Meaning \\
\midrule
$N$ & Number of neurons per region-session ($N=1{,}000$). \\
$x_i(t)$ & Deconvolved activity of neuron $i \in \{1,\dots,N\}$ at time $t$. \\
$[w_0, w_1)$ & Analysis window on the session clock. \\
$\Delta_b$ & Temporal bin width ($1.25$~s for dimensionality; none, i.e.\ native frames, for prediction). \\
$T$ & Number of time points (bins or frames) in the window; $t_1,\dots,t_T$ their times. \\
$(m_i, a_i)$ & ML and AP coordinate of neuron $i$ in \textmu m. \\
$\delta$ & Checkerboard square size, $\delta = 60$~\textmu m. \\
$\mathcal{A}, \mathcal{B}$ & The two spatially interleaved neuron sets; $N_A = |\mathcal{A}|$, $N_B = |\mathcal{B}|$. \\
$F \in \R^{N_A \times T}$, $G \in \R^{N_B \times T}$ & Mean-centred activity of sets $\mathcal{A}$ and $\mathcal{B}$. \\
$L$, $\epsilon$ & Train/test block length ($L = 72$~s) and padding removed at each block edge ($\epsilon = 2$~s). \\
$\mathcal{T}_{\mathrm{tr}}, \mathcal{T}_{\mathrm{te}}$ & Training and test time points; $T_{\mathrm{tr}}, T_{\mathrm{te}}$ their numbers. \\
$C \in \R^{N_A \times N_B}$ & Training cross-covariance between the two sets. \\
$K$ & Number of shared variance components (SVCs), $K = \min(K_{\max}, N_A, N_B, T_{\mathrm{tr}}, T_{\mathrm{te}})$, $K_{\max}=128$. \\
$\mathbf{u}_k, \mathbf{v}_k$ & $k$-th left/right singular vector of $C$ (SVC $k$ in sets $\mathcal{A}$/$\mathcal{B}$). \\
$p_k(t), q_k(t)$ & Projections of the two sets onto SVC $k$. \\
$\hat{s}_k$ & Reliable (cross-set) variance of SVC $k$ on test data. \\
$\sigma_k^2$ & Total variance of SVC $k$ on test data (mean of both sets). \\
$r_k$ & Reliable-variance fraction (``maximum explainable'') of SVC $k$. \\
$\alpha$ & Power-law decay exponent of $r_k$. \\
$d_{50}$ & Effective dimensionality (number of SVCs holding 50\% of reliable variance). \\
$Z \in \R^{M \times T}$ & Predictor matrix with $M$ predictors (video PCs, wheel speed, block or choice). \\
$e^{A}_k(t), e^{B}_k(t)$ & Test-set residuals of $p_k$, $q_k$ after regression on $Z$. \\
$v_k$ & Fraction of the variance of SVC $k$ explained by $Z$. \\
$n$ & Number of sessions averaged in a region. \\
\bottomrule
\end{longtable}

\subsection{Shared variance component analysis (SVCA)}

SVCA follows Stringer et al.~\cite{stringer2019}: it estimates, dimension by dimension, how much of the population's variance is shared between two non-overlapping sets of neurons and therefore reproducible rather than independent noise.

\paragraph{Binning.} For the dimensionality analysis the window is divided into consecutive bins $[w_0 + j\Delta_b,\, w_0 + (j+1)\Delta_b)$ and each neuron's activity is averaged over the frames falling into each bin; bins without frames are dropped. The bin centre is the time of each bin.

\paragraph{Spatial split.} To prevent neighbouring neurons (which may share out-of-focus fluorescence) from ending up in different sets and inflating the shared variance, neurons are assigned to a two-dimensional checkerboard in the ML--AP plane:
\begin{equation}
  \xi_i = \mathbb{1}\!\left[m_i \bmod 2\delta < \delta\right], \quad
  \eta_i = \mathbb{1}\!\left[a_i \bmod 2\delta < \delta\right], \quad
  i \in \mathcal{A} \iff (\xi_i + \eta_i) \bmod 2 = 0,
\end{equation}
and $\mathcal{B}$ contains all other neurons. Each row of $F$ and $G$ is centred by subtracting its mean over the whole window.

\paragraph{Train/test partition.} Time points are assigned to alternating contiguous blocks of length $L$, with block index $b(t) = \lfloor (t - t_1)/L \rfloor$ (where $t_1$ is the first time point):
\begin{equation}
  t \in \mathcal{T}_{\mathrm{tr}} \iff b(t) \text{ even}, \qquad t \in \mathcal{T}_{\mathrm{te}} \iff b(t) \text{ odd},
\end{equation}
and every time point within $\epsilon$ of a block boundary, i.e.\ with $(t - t_1) \bmod L < \epsilon$ or $> L - \epsilon$, is removed from both sets to limit leakage through slow autocorrelation.

\paragraph{Shared components.} On training data the cross-covariance
\begin{equation}
  C = \frac{1}{T_{\mathrm{tr}}} F_{\mathrm{tr}} G_{\mathrm{tr}}^{\top}
\end{equation}
is decomposed by a truncated (randomised) singular value decomposition, $C \approx \sum_{k=1}^{K} s_k \mathbf{u}_k \mathbf{v}_k^{\top}$ with $s_1 \ge s_2 \ge \dots$. On test data the two sets are projected onto their own singular vectors,
\begin{equation}
  p_k(t) = \mathbf{u}_k^{\top} F_{:,t}, \qquad q_k(t) = \mathbf{v}_k^{\top} G_{:,t}, \qquad t \in \mathcal{T}_{\mathrm{te}},
\end{equation}
and the reliable variance, the total variance and their ratio are
\begin{equation}
  \hat{s}_k = \frac{1}{T_{\mathrm{te}}} \sum_{t \in \mathcal{T}_{\mathrm{te}}} p_k(t)\, q_k(t), \qquad
  \sigma_k^2 = \frac{1}{2 T_{\mathrm{te}}} \sum_{t \in \mathcal{T}_{\mathrm{te}}} \left( p_k(t)^2 + q_k(t)^2 \right), \qquad
  r_k = \frac{\hat{s}_k}{\sigma_k^2}.
\end{equation}
Because noise independent between the two sets does not contribute to $\hat{s}_k$ in expectation, $r_k$ estimates the fraction of SVC $k$'s variance that is shared and hence the maximum that any predictor can explain. On finite data $r_k$ can be negative for weak components.

\paragraph{Power-law exponent.} The decay of the spectrum is summarised by fitting $|r_k| \approx c\, k^{-\alpha}$ over ranks $k \in \mathcal{K} = \{10, \dots, 100\}$ by weighted least squares in log--log coordinates, with weights $w_k = 1/k$ so that the noisier high ranks do not dominate (adapted from MouseLand \texttt{critical\_init}~\cite{criticalinit}):
\begin{equation}
  (\hat{\alpha}, \hat{c}) = \arg\min_{\alpha, c} \sum_{k \in \mathcal{K}} w_k \left( \log |r_k| + \alpha \log k - \log c \right)^2 .
\end{equation}

\paragraph{Effective dimensionality.} With $r_k^{+} = \max(r_k, 0)$,
\begin{equation}
  d_{50} = \min \Big\{ K' : \textstyle\sum_{k=1}^{K'} r_k^{+} \ge 0.5 \sum_{k=1}^{K} r_k^{+} \Big\},
\end{equation}
i.e.\ the number of leading SVCs holding half of the captured reliable variance. This measure is independent of the absolute magnitude of $r_1$.

\subsection{Predicting shared variance from behavioural and task variables}

\paragraph{Predictors.} All predictors are brought to the neural time points $t_1, \dots, t_T$ (native imaging frames in this analysis).
\begin{itemize}
  \item \textit{Left-camera video motion PCs} ($M = 16$). Left-camera frames $I_\tau$ within the window are converted to greyscale and downsampled to $60 \times 45$ pixels (area averaging), and motion energy is $\mathbf{m}_\tau = |I_\tau - I_{\tau-1}| \in \R^{2700}$. Following Stringer et al.~\cite{stringer2019}, a two-stage SVD is used: the column-centred motion-energy matrix of each 300-s video segment (here the window is a single segment) is decomposed, and its top 200 spatial singular vectors, scaled by their singular values, are kept. These spatial summaries are concatenated across segments and decomposed again to obtain 16 spatial components $\mathbf{w}_1, \dots, \mathbf{w}_{16}$. The PCs are the projections $\mathbf{w}_j^{\top} \mathbf{m}_\tau$, each assigned to the nearest neural time point.
  \item \textit{Wheel speed} ($M = 1$). Wheel position is linearly interpolated at 250~Hz, low-pass filtered (8th-order Butterworth, 20~Hz corner, zero-phase), and differentiated. The absolute value is assigned to the nearest neural time point.
  \item \textit{Block} ($M = 1$). The block prior $\pi_{\mathrm{L}} \in \{0.2, 0.5, 0.8\}$ of the trial in progress, held constant from each trial's start until the next trial's start.
  \item \textit{Choice} ($M = 1$). $c \in \{-1, +1\}$ during the 0.5~s following first-movement onset of each trial and $0$ otherwise; no-go trials and trials without a detected movement are excluded.
\end{itemize}
Whisker-pad motion energy is not included in the wheel predictor because it is not available for all sessions. In the passive protocol there are no trials, and extracted wheel data cover the passive epoch in only two sessions, so the passive behavioural analysis uses the left-camera video PCs only (computed from the 730-s passive window, i.e.\ three video segments).

\paragraph{Movement during the passive protocol.} To test whether the animals are still during the passive protocol, the mean whisker-pad motion energy $\bar{E}$ (IBL \texttt{leftCamera.ROIMotionEnergy}) is compared between the passive window and the duration-matched task window, $\rho = \bar{E}_{\mathrm{passive}} / \bar{E}_{\mathrm{task}}$. This is possible for the ten region-sessions in which motion energy was extracted.

\paragraph{Regression.} Each predictor is z-scored with its training-set mean and standard deviation; predictors that are constant on the training set are dropped, and a predictor set with no varying predictor explains nothing ($v_k = 0$; this happens for the block predictor when the 300-s window lies inside a single block). The SVC projections of each set are regressed on $Z$ independently, by unregularised least squares on training data,
\begin{equation}
  B^{A} = \arg\min_{B} \sum_{t \in \mathcal{T}_{\mathrm{tr}}} \big\| \mathbf{p}(t) - B\, \mathbf{z}(t) \big\|^2, \qquad
  B^{B} = \arg\min_{B} \sum_{t \in \mathcal{T}_{\mathrm{tr}}} \big\| \mathbf{q}(t) - B\, \mathbf{z}(t) \big\|^2,
\end{equation}
with $\mathbf{p}(t) = (p_1(t), \dots, p_K(t))^{\top}$, $\mathbf{q}(t)$ analogously and $\mathbf{z}(t) = Z_{:,t}$. The test-set residuals are $e^{A}_k(t) = p_k(t) - [B^{A} \mathbf{z}(t)]_k$ and $e^{B}_k(t) = q_k(t) - [B^{B} \mathbf{z}(t)]_k$. The shared variance that survives the prediction is their cross-covariance, and the explained fraction is
\begin{equation}
  v_k = \frac{\hat{s}_k - \frac{1}{T_{\mathrm{te}}} \sum_{t \in \mathcal{T}_{\mathrm{te}}} e^{A}_k(t)\, e^{B}_k(t)}{\sigma_k^2},
\end{equation}
so that $v_k \le r_k$ when prediction removes only shared variance. Because the two sets are regressed separately, only variance that the predictor explains in both sets reduces the cross-covariance. For components without reliable shared variance ($r_k \lesssim 0$, typically at high ranks) the test-set residual cross-covariance can become negative, so that $v_k$ can exceed $r_k$; $v_k$ is only interpretable where $r_k$ is clearly positive.

\subsection{Averaging across sessions}

For each region, per-session curves are averaged component by component over the $n$ sessions, $\bar{y}_k = \frac{1}{n}\sum_{j=1}^{n} y_{jk}$, and the band shows the standard error $\mathrm{SEM}_k = \mathrm{sd}_k / \sqrt{n}$. In Fig.~\ref{fig:dim}a the $r_k$ are averaged as they are and $\mathrm{sd}_k$ is the population standard deviation (normalised by $n$). In Fig.~\ref{fig:behav} each session's $r_k$ and $v_k$ are clipped at zero before averaging and $\mathrm{sd}_k$ is the sample standard deviation (normalised by $n - 1$). Because sessions from the same animal are not independent, and MOp and MOs come from few animals, this uncertainty is between-session, not between-animal.

\subsection{Software and code availability}

Analyses were run in Python (IBL \texttt{ONE}, \texttt{iblatlas}, \texttt{brainbox}, NumPy, SciPy, OpenCV); truncated SVDs use \texttt{torch.svd\_lowrank} on the local GPU. The implementation lives in the IBL mesoscope repository, \url{https://github.com/int-brain-lab/mesoscope}~\cite{mesoscoperepo}: the data loader (\path{meso_loader.py}), SVCA and behavioural prediction (\path{pilot_analyses/stringer19_svca_prediction.py}), and regional neuron selection, power-law fit and effective dimensionality (\path{pilot_analyses/stringer19_region_comparison.py}). The scripts that select the sessions, produce each figure and generate the tables of this report are in \path{pilot_analyses/stringer2019/} (\path{region_session_scan.py}, \path{regional_dimensionality.py}, \path{regional_behavior_prediction.py}, \path{passive_regional_dimensionality.py}, \path{passive_behavior_prediction.py}, \path{make_tables.py}); its \path{README.md} gives the order in which to run them.

\section{Results}

\subsection{Task epoch}

\begin{figure}[p]
  \centering
  \includegraphics[width=\textwidth]{../output/regional_dimensionality_panels.pdf}
  \caption{\textbf{Regional comparison of shared-variance dimensionality.}
  \textbf{a} Mean reliable-variance spectra $r_k$ across sessions for VISp, SSp-bfd, MOp and MOs; shading denotes the session-level SEM ($n = 5$ sessions per region; VISp and SSp-bfd from four animals each, MOp from two, MOs from one).
  \textbf{b} Power-law decay exponent $\alpha$ fitted over SVC ranks 10--100; dots are sessions, bars means.
  \textbf{c} Effective dimensionality $d_{50}$.
  Each region-session: $N = 1{,}000$ neurons, 100--3,100~s, 1.25-s bins, $K_{\max} = 128$.}
  \label{fig:dim}
\end{figure}

\begin{figure}[p]
  \centering
  \includegraphics[width=\textwidth]{../output/regional_behavior_prediction_sessions_panels.pdf}
  \caption{\textbf{Behavioural and task prediction of shared neural variance across cortical regions.}
  \textbf{a--d} Component-wise maximum explainable variance $r_k$ and variance $v_k$ explained by left-camera video motion PCs, wheel speed, block or choice in \textbf{a} MOp, \textbf{b} MOs, \textbf{c} VISp and \textbf{d} SSp-bfd. Thick lines: mean across sessions; shading: SEM; thin lines: the $n = 8$ individual sessions per region (VISp and SSp-bfd from eight animals each, MOp from three, MOs from one). Each region-session: $N = 1{,}000$ neurons, 1,730--2,030~s at native frame rate, $K_{\max} = 128$.}
  \label{fig:behav}
\end{figure}

\paragraph{Dimensionality (Fig.~\ref{fig:dim}).} The reliable-variance spectra of the four regions are similar in magnitude (mean $r_1$: VISp 45.1\%, SSp-bfd 43.1\%, MOp 45.8\%, MOs 44.8\%) and decay over similar ranges. VISp has the steepest decay ($\alpha = 0.81 \pm 0.14$, mean $\pm$ s.d.\ across sessions) and SSp-bfd the shallowest ($0.58 \pm 0.17$), with MOp ($0.72 \pm 0.13$) and MOs ($0.73 \pm 0.14$) in between. Effective dimensionality is 26--30 components in all regions (VISp $28.8 \pm 3.4$, SSp-bfd $29.6 \pm 2.7$, MOp $26.6 \pm 3.4$, MOs $25.6 \pm 6.0$). With five sessions per region, neither measure differs significantly between regions (Kruskal--Wallis test across the four regions: $\alpha$, $p = 0.13$; $d_{50}$, $p = 0.46$).

\paragraph{Behavioural prediction (Fig.~\ref{fig:behav}).} The individual sessions show that the SEM bands of the earlier three-session version come from genuine between-session variability: the maximum explainable variance of SVC~1 ranges over tens of percentage points between sessions of the same region. Averaged over eight sessions, $r_1$ is 29.0 $\pm$ 3.5\% in MOp, 17.3 $\pm$ 3.6\% in MOs, 25.9 $\pm$ 5.3\% in VISp and 21.3 $\pm$ 3.2\% in SSp-bfd (mean $\pm$ SEM). The lower values than in Fig.~\ref{fig:dim} are expected from the shorter window and the native (unbinned) frame rate. Only a small part of the reliable variance is predicted by behaviour. Summed over the ten leading SVCs, left-camera video motion PCs explain on average 17\% of the reliable variance in MOp, 14\% in MOs and SSp-bfd, and 5\% in VISp; wheel speed explains 10\%, 8\%, 10\% and 5\%, respectively. Block and choice each explain at most 2.5\%. The block predictor is uninformative in seven of the 32 region-sessions, because the 300-s window lies within a single block. Per-session values for the left-camera video span 0--27\% of the reliable variance, so the regional differences in behavioural predictability are small relative to the variability between sessions.

\subsection{Passive protocol}

\begin{figure}[p]
  \centering
  \includegraphics[width=\textwidth]{../output/passive_regional_dimensionality_panels.pdf}
  \caption{\textbf{Regional shared-variance dimensionality during the passive protocol.}
  \textbf{a} Mean reliable-variance spectra $r_k$ from the first 730~s of imaged passive protocol; shading: session-level SEM ($n = 5$ sessions per region; VISp and SSp-bfd from four animals each, MOp from three, MOs from one).
  \textbf{b} Power-law exponent $\alpha$ and \textbf{c} effective dimensionality $d_{50}$ in the passive window (filled) and in a duration-matched 730-s task window (100--830~s; open) from the same neurons; lines connect sessions, bars mark means.
  \textbf{d} Whisker-pad motion energy in the passive window relative to the task window, $\rho$, for the ten region-sessions with extracted motion energy; dashed line: equal movement.
  Each region-session: $N = 1{,}000$ neurons, 1.25-s bins, $K_{\max} = 128$.}
  \label{fig:pdim}
\end{figure}

\begin{figure}[p]
  \centering
  \includegraphics[width=\textwidth]{../output/passive_behavior_prediction_panels.pdf}
  \caption{\textbf{Left-camera video prediction of shared neural variance during the passive protocol.}
  \textbf{a--d} Maximum explainable variance $r_k$ and variance $v_k$ explained by left-camera video motion PCs in \textbf{a} MOp, \textbf{b} MOs, \textbf{c} VISp and \textbf{d} SSp-bfd, from the first 730~s of imaged passive protocol. Thick lines: mean; shading: SEM; thin lines: the $n = 5$ individual sessions per region. Each region-session: $N = 1{,}000$ neurons, native frame rate, $K_{\max} = 128$, 16 video PCs. The large video values at SVC~19--80 in one VISp session (\texttt{76f74094}) lie where its $r_k$ is near or below zero, where $v_k$ is not interpretable.}
  \label{fig:pbehav}
\end{figure}

\paragraph{Dimensionality (Fig.~\ref{fig:pdim}).} Restricting the analysis to the passive protocol does not change the regional picture qualitatively. Compared with a duration-matched task window from the same neurons, neither the power-law exponent (passive $0.32$ vs.\ task $0.44$, mean over the 20 region-sessions; Wilcoxon signed-rank $p = 0.11$) nor the effective dimensionality ($39.2$ vs.\ $35.5$; $p = 0.12$) differs significantly. The numerically lower passive $\alpha$ comes from MOp ($0.35$ vs.\ $0.61$) and MOs ($0.16$ vs.\ $0.37$), not from VISp ($0.61$ vs.\ $0.54$) or SSp-bfd ($0.16$ vs.\ $0.24$). The 730-s windows give noisier spectra and lower $\alpha$ than the 3,000-s analysis of Fig.~\ref{fig:dim}, so absolute values should only be compared within Fig.~\ref{fig:pdim}. Importantly, the passive protocol is not a movement-free epoch: whisker-pad motion energy during passive is 60--95\% of that in the task window (mean $\rho = 0.79$).

\paragraph{Behavioural prediction (Fig.~\ref{fig:pbehav}).} Consistent with the continued movement, the left-camera video still predicts a substantial part of the shared variance during the passive protocol, most clearly in the motor areas. Summed over the ten leading SVCs, video PCs explain on average 25\% of the reliable variance in MOs (range across sessions 17--34\%), 21\% in MOp (0--36\%), 11\% in SSp-bfd (1--18\%) and 5\% in VISp (1--7\%). For SVC~1, video explains $13.9 \pm 2.6\%$ (MOs), $11.9 \pm 4.3\%$ (MOp), $7.5 \pm 3.0\%$ (SSp-bfd) and $1.1 \pm 0.7\%$ (VISp) of the variance (mean $\pm$ SEM), out of a maximum explainable $28.2$, $24.7$, $28.0$ and $29.6\%$. The ordering of regions is similar to the task epoch (Fig.~\ref{fig:behav}), with VISp the least predictable from the left-camera video. The windows differ in duration and position, however (730~s of passive protocol vs.\ 300~s of task), and the session sets only partly overlap, so the two figures should not be compared quantitatively.

\subsection{Caveats}
 (i) Uncertainty is between sessions, not between animals; MOs is represented by a single animal and MOp by two (Fig.~\ref{fig:dim}) or three (Fig.~\ref{fig:behav}), so regional differences involving the motor areas cannot be separated from animal identity. (ii) The behavioural analysis uses a single 300-s window per session, a 16-dimensional video representation and unregularised regression, so its prediction estimates are noisy and conservative relative to Stringer et al.~\cite{stringer2019}, who used longer recordings and 1,000 video PCs. (iii) Wheel speed is the only non-video behavioural predictor; whisker-pad motion energy is not available for every session.

\begin{thebibliography}{9}
\bibitem{stringer2019} C.~Stringer, M.~Pachitariu, N.~Steinmetz, C.~B.~Reddy, M.~Carandini and K.~D.~Harris. Spontaneous behaviors drive multidimensional, brainwide activity. \textit{Science} \textbf{364}, eaav7893 (2019). \href{https://doi.org/10.1126/science.aav7893}{doi:10.1126/science.aav7893}.
\bibitem{criticalinit} MouseLand. \texttt{critical\_init}. \url{https://github.com/MouseLand/critical_init}.
\bibitem{mesoscoperepo} International Brain Laboratory. \texttt{mesoscope}: loader and analyses for IBL mesoscope recordings. \url{https://github.com/int-brain-lab/mesoscope}.
\end{thebibliography}

\clearpage
\appendix
\section{Sessions}

\input{../output/table_dimensionality.tex}

\input{../output/table_behavior.tex}

\input{../output/table_passive.tex}

\end{document}
